\documentclass[11pt]{article}
% ===========================================================================
%  Shared preamble for the AP Calculus AB exam-preparation worksheets.
%  Input right after \documentclass. Files beginning with "_" are not compiled
%  on their own. Navy and gold, to match the course's lesson pages.
% ===========================================================================
\usepackage[letterpaper,top=2.3cm,bottom=2.2cm,left=1.7cm,right=1.7cm,headheight=28pt,footskip=30pt]{geometry}
\usepackage{amsmath,amssymb}
\usepackage[T1]{fontenc}
\usepackage[utf8]{inputenc}
\usepackage{lmodern}
\usepackage{xcolor}
\usepackage[most]{tcolorbox}
\usepackage{enumitem}
\usepackage{fancyhdr}
\usepackage{array}
\usepackage{tikz}
\usepackage{pgfplots}
\pgfplotsset{compat=1.18}

\definecolor{navy}{HTML}{161E54}
\definecolor{ink}{HTML}{27357E}
\definecolor{gold}{HTML}{B98A00}
\definecolor{goldbg}{HTML}{FFFDF3}
\definecolor{inkbg}{HTML}{F6F7FD}
\definecolor{rule}{HTML}{C9CFF0}
\definecolor{okgreen}{HTML}{1B7A44}
\definecolor{okbg}{HTML}{F1F8F3}

\pagestyle{fancy}
\fancyhf{}
\fancyhead[L]{\footnotesize Name: \underline{\hspace{4cm}}}
\fancyhead[C]{\textbf{\color{navy}Integration Academy}\\[-2pt]{\scriptsize Calculus AB \textperiodcentered{} AP\textsuperscript{\textregistered} Exam Preparation}}
\fancyfoot[L]{\tiny\color{black!55}\parbox[t]{0.78\textwidth}{AP\textsuperscript{\textregistered} is a registered trademark of the College Board, which was not involved in the production of, and does not endorse, this material. \textcopyright{} Integration Academy.}}
\fancyfoot[R]{\scriptsize\thepage}
\renewcommand{\headrulewidth}{1.2pt}
\renewcommand{\headrule}{\hbox to\headwidth{\color{gold}\leaders\hrule height \headrulewidth\hfill}}
\renewcommand{\footrulewidth}{0pt}
\setlength{\parindent}{0pt}
\setlength{\parskip}{4pt}

% Title block: \wstitle{code}{title}{unit line}
\newcommand{\wstitle}[3]{\begin{center}{\small\bfseries\color{gold}WORKSHEET #1}\\[3pt]{\LARGE\bfseries\color{navy}#2}\\[3pt]{\small\color{black!60}#3}\end{center}\vspace{2pt}}
% Section heading
\newcommand{\sech}[1]{\par\vspace{10pt}{\large\bfseries\color{navy}#1}\par\nobreak\vspace{2pt}{\color{rule}\hrule height0.8pt}\vspace{7pt}}
% A level tag printed beside a group of questions
\newcommand{\level}[1]{\par\vspace{6pt}{\footnotesize\bfseries\color{gold}\MakeUppercase{#1}}\par\nobreak\vspace{2pt}}

\newtcolorbox{ideabox}{colframe=ink,colback=inkbg,coltitle=white,colbacktitle=ink,fonttitle=\bfseries\small,title={Key ideas},arc=2pt,boxrule=0.8pt,left=8pt,right=8pt}
\newtcolorbox{example}[1]{colframe=rule,colback=white,coltitle=navy,colbacktitle=inkbg,fonttitle=\bfseries\small,title={#1},arc=2pt,boxrule=0.8pt,borderline west={3pt}{0pt}{ink},left=10pt,right=8pt}
\newtcolorbox{frbox}[1]{colframe=gold,colback=goldbg,coltitle=white,colbacktitle=gold,fonttitle=\bfseries\small,title={#1},arc=2pt,boxrule=0.8pt,breakable,left=8pt,right=8pt}
\newtcolorbox{solbox}{colframe=okgreen,colback=okbg,coltitle=white,colbacktitle=okgreen,fonttitle=\bfseries\small,title={Worked solutions},arc=2pt,boxrule=0.8pt,breakable,left=8pt,right=8pt}
\newtcolorbox{scorebox}[1]{colframe=navy,colback=white,coltitle=white,colbacktitle=navy,fonttitle=\bfseries\small,title={#1},arc=2pt,boxrule=0.8pt,breakable,left=8pt,right=8pt}

\newcommand{\soln}{\par\smallskip\textbf{\color{ink}Solution.}\ }
\newcommand{\result}[1]{\par\smallskip\textbf{#1}}
% blank working space of a given height
\newcommand{\work}[1]{\par\vspace{#1}}
% a point value in the margin of a scoring guide
\newcommand{\pts}[1]{\hfill{\footnotesize\bfseries\color{gold}[#1]}}

% Coordinate grid for graph questions: \begin{apgraph}{xmin}{xmax}{ymin}{ymax} ... \end{apgraph}
\newenvironment{apgraph}[4]{\begin{center}\begin{tikzpicture}\begin{axis}[
  axis lines=middle, xmin=#1, xmax=#2, ymin=#3, ymax=#4,
  grid=both, grid style={black!12}, xtick distance=1, ytick distance=1,
  ticklabel style={font=\scriptsize}, xlabel={$x$}, ylabel={$y$},
  width=10cm, height=7.2cm, samples=80, clip=false,
  every axis plot/.append style={line width=1.1pt, color=ink}]}{\end{axis}\end{tikzpicture}\end{center}}
\newcommand{\closeddot}[2]{\addplot[only marks, mark=*, mark size=2.2pt, color=ink] coordinates {(#1,#2)};}
\newcommand{\opendot}[2]{\addplot[only marks, mark=*, mark size=2.2pt, color=ink, mark options={fill=white}] coordinates {(#1,#2)};}

% Multiple-choice options: four across, or stacked when the options are long.
\newcommand{\opts}[4]{\par\vspace{3pt}\begin{tabular}{@{}p{0.235\linewidth}p{0.235\linewidth}p{0.235\linewidth}p{0.235\linewidth}@{}}(A)\ #1 & (B)\ #2 & (C)\ #3 & (D)\ #4\end{tabular}\par\vspace{9pt}}
\newcommand{\optsv}[4]{\par\vspace{3pt}\begin{tabular}{@{}l@{\ \ }p{0.9\linewidth}@{}}(A) & #1\\(B) & #2\\(C) & #3\\(D) & #4\end{tabular}\par\vspace{9pt}}

\fancyhead[R]{\footnotesize\color{black!60}Worksheet 9.5}
\begin{document}

\wstitle{9.5}{Free-Response Type: Functions Given by Graphs and Tables}{Unit 9: Exam Strategy and Free-Response Types}

\begin{ideabox}
\begin{itemize}[leftmargin=*,itemsep=2pt,topsep=2pt]
  \item From a graph of $f'$: $f$ increases where $f'>0$; $f$ is concave up where $f'$ is increasing; $f''$ is the slope of the graph.
  \item Values of $f$ come from areas: $f(b)=f(a)+\displaystyle\int_{a}^{b}f'(x)\,dx$, with areas below the axis counted as negative.
  \item From a table: estimate a derivative with the two nearest points, and an integral with a Riemann or trapezoidal sum.
  \item $\displaystyle\int_{a}^{b}f'(x)\,dx=f(b)-f(a)$ exactly, whatever the spacing of the table.
\end{itemize}
\end{ideabox}

\sech{Worked examples}
\begin{example}{Example 1 \textperiodcentered{} A relative maximum}
The three examples use the graph of $f'$ printed at the top of the next page, with $f(-1)=3$. At what value of $x$ does $f$ have a relative maximum?\soln $f'$ changes from positive to negative where the graph crosses the axis going down.\result{$x=0$.}
\end{example}

\begin{example}{Example 2 \textperiodcentered{} A value from areas}
Find $f(5)$.\soln $f(5)=3+\displaystyle\int_{-1}^{5}f'(x)\,dx=3+(1-1-2+4)$.\result{$f(5)=5$.}
\end{example}

\begin{example}{Example 3 \textperiodcentered{} Concavity}
On what interval is the graph of $f$ concave up?\soln $f'$ is increasing on $(1,3)$ and on $(3,5)$.\result{$1<x<5$.}
\end{example}

\newpage
\sech{Practice}
\emph{No calculator on this page.} The graph of $f'$, the derivative of $f$, consists of four line segments, as shown. It is given that $f(-1)=3$. The worked examples on the first page and Questions 1 to 4 and 9 use this graph.
\begin{apgraph}{-4}{6}{-3}{5}
\addplot[color=ink, mark=none, line width=1.1pt] coordinates {(-3,1) (-1,2) (1,-2) (3,0) (5,4)};
\closeddot{-3}{1}\closeddot{5}{4}
\node[font=\small, color=ink] at (axis cs:3.2,2.6) {$y=f'(x)$};
\end{apgraph}
\level{Core}
\begin{enumerate}[leftmargin=*,itemsep=2pt]
  \item Find $f'(-2)$ and $f''(-2)$.\work{1.1cm}
  \item Find $f''(4)$.\work{1.1cm}
  \item Find $f(-3)$.\work{1.1cm}
  \item On what interval is $f$ decreasing? Give a reason.\work{1.1cm}
\end{enumerate}

\level{Exam level}
\begin{enumerate}[leftmargin=*,itemsep=2pt,start=5]
  \item The table gives the temperature $H(t)$, in $^\circ$C, of a liquid at selected times $t$ minutes. Estimate $H'(5)$, with units.\begin{center}\renewcommand{\arraystretch}{1.25}\begin{tabular}{|c|c|c|c|c|c|}\hline $t$ & $0$ & $4$ & $6$ & $10$ & $14$ \\ \hline $H(t)$ & $20$ & $32$ & $38$ & $44$ & $46$ \\ \hline\end{tabular}\end{center}\work{2.1cm}
  \item Use a trapezoidal sum with the four subintervals of the table to estimate $\displaystyle\int_{0}^{14}H(t)\,dt$.\work{2.1cm}
  \item Use your answer to Question 6 to estimate the average temperature for $0\le t\le14$.\work{2.1cm}
  \item Evaluate $\displaystyle\int_{0}^{14}H'(t)\,dt$ and state its meaning.\work{2.1cm}
\end{enumerate}

\begin{frbox}{Question 9 \textperiodcentered{} Free-response practice}
Use the graph of $f'$ at the top of the page, with $f(-1)=3$.
\begin{enumerate}[label=(\alph*),leftmargin=*,itemsep=2pt]
  \item Find the absolute minimum value of $f$ on $[-3,5]$. Justify your answer.\work{1.9cm}
  \item Find the $x$-coordinate of each point of inflection of the graph of $f$. Give a reason.\work{1.9cm}
  \item Let $h(x)=x\,f(x)$. Find $h'(1)$.\work{1.9cm}
\end{enumerate}
\end{frbox}

\newpage
\begin{solbox}
\begin{enumerate}[leftmargin=*,itemsep=4pt]
  \item $f'(-2)=\dfrac32$ and $f''(-2)=\dfrac12$.
  \item $f''(4)=2$.
  \item $f(-3)=3-\displaystyle\int_{-3}^{-1}f'(x)\,dx=3-3=0$.
  \item $0<x<3$, because $f'(x)<0$ there.
  \item $\dfrac{38-32}{6-4}=3$ $^\circ$C per minute.
  \item $\dfrac{4(52)}{2}+\dfrac{2(70)}{2}+\dfrac{4(82)}{2}+\dfrac{4(90)}{2}=104+70+164+180=518$.
  \item $\dfrac{518}{14}=37$ $^\circ$C.
  \item $H(14)-H(0)=26$: the temperature rises 26 $^\circ$C over the 14 minutes.
  \item \begin{enumerate}[label=(\alph*),leftmargin=*,itemsep=3pt]
    \item Candidates are the endpoints and $x=3$, where $f'$ changes from negative to positive. $f(-3)=0$, $f(3)=3+1-1-2=1$ and $f(5)=5$. The absolute minimum value is 0, at $x=-3$.
    \item $x=-1$ and $x=1$: $f'$ changes from increasing to decreasing at $x=-1$ and from decreasing to increasing at $x=1$.
    \item $h'(x)=f(x)+x\,f'(x)$, and $f(1)=3+1-1=3$, $f'(1)=-2$: $h'(1)=3-2=1$.
  \end{enumerate}
\end{enumerate}
\end{solbox}

\end{document}
